\begin{abstract}
\noindent
Large language models may seem hard to beat at extracting information from documents. But using an external API for every document can become costly. We wanted to find out whether smaller models adapted to a specific task could achieve comparable results locally. Magda turns flyers from the German supermarket chain PENNY into structured offers by finding product details and linking those that belong together. For our Information Extraction seminar, we followed our project proposal by using Claude Sonnet 5 to label 666 pages with 33,457 entity spans. We used the training split to fine-tune GBERT, XLM-R, LiLT and LayoutXLM on remote GPUs, comparing models using text, word positions and page images. We also trained a multilayer perceptron to predict whether two entities belong to the same offer based on entity types, spatial relations and lexical features. We compare this learned grouping approach with a rule-based heuristic and evaluate Union-Find and integer linear programming for constructing offer groups. On 42 pages with human annotations, Magda matches 212 of 279 reference records among 281 predictions and achieves an F1 score of 0.7571 for product names and regular prices. Qwen3.8-27B returns 349 records, including some duplicates, matches 240 reference records and achieves 0.7643 F1. The 95\% paired bootstrap interval for the F1 difference includes zero, providing no clear evidence that either system performs better. Entity recognition is evaluated separately using SemEval matching schemes to distinguish exact span matches from partial overlaps and type agreement. Record F1 measures agreement on product names and regular prices rather than the correctness of complete offers. Known annotation errors and prior inspection of test pages limit the conclusions to an exploratory comparison.
\end{abstract}
\begin{quote}\small
\textit{Keywords:} information extraction, document understanding, supermarket flyers, entity recognition, relation grouping, weak supervision.
\end{quote}
